\chapter[Pojemność Sobolewa i osobliwości]{Pojemność Sobolewa, osobliwości usuwalne i regularność brzegu}
\label{chap:pde19}

\section{Dlaczego miara Lebesgue'a nie wystarcza?}
\label{sec:pde19-motivation}
W teorii słabych rozwiązań PDE rozróżniamy zbiory pomijalne w sensie
miary Lebesgue'a od zbiorów pomijalnych dla energii Dirichleta.
Przykładowo pojedynczy punkt w $\R^2$ ma miarę zero, a funkcja
$\log|x|$ jest harmoniczna poza tym punktem, gdyż
dla $r=|x|>0$ mamy $\Delta\log r=-r^{-2}+r^{-2}=0$.
Jej gradient nie jest jednak całkowalny z kwadratem
w otoczeniu zera:
\[
 \int_{B_R(0)}|\nabla\log|x||^2\,\mathrm d\lambda_2(x)
 =2\pi\int_0^R\frac{1}{r}\,\mathrm dr=+\infty.
\] Tymczasem funkcja
harmoniczna, która \emph{należy do $H^1$ także w pobliżu osobliwości},
może nie odczuwać usunięcia punktu. Sformułujemy tę zasadę przy
pomocy pojemności.

Zawsze, gdy piszemy $\Omega\subset\R^d$ jako \emph{obszar},
zakładamy, że jest on otwarty i spójny. W definicjach lokalnych
pojemności nie wymagamy spójności zbioru otwartego.
Rozpatrujemy funkcje rzeczywiste; argumenty dotyczące
równań z rzeczywistymi współczynnikami dla funkcji zespolonych
prowadzi się osobno dla części rzeczywistej i urojonej.
Przyjmujemy $|E|=\lambda_d(E)$ dla zbiorów mierzalnych,
$\omega_d=|B_1(0)|$ i $\|\cdot\|_p=\|\cdot\|_{L^p(\Omega)}$,
o ile zbiór, na którym liczymy normę, nie został wskazany inaczej.

\begin{uwaga}[Równość prawie wszędzie a quasi-ciągłość]
\label{rem:pde-v43-capacity-ae}
Element $H^1(\Omega)$ jest klasą funkcji równych prawie
wszędzie; w tym sensie nie ma zadanego punktowego śladu
na dowolnym zbiorze miary zero. Pojemność oparta na energii
pozwala rozróżniać zbiory, na których reprezentant
quasi-ciągły może mieć znaczenie. Nie wynika z tego, że
każdy zbiór miary Lebesgue'a zero ma zerową pojemność.
Relację z operatorem śladu na brzegu Lipschitzowskim
omówiono w rozdziale~\ref{chap:pde-embeddings}; są to
różne konstrukcje, wymagające odmiennych założeń.
\end{uwaga}

\section{Definicja pojemności i jej podstawowe własności}
\label{sec:pde19-capacity}
\begin{definicja}[Względna pojemność Dirichleta]
\label{def:pde19-cap}
Niech $\Omega\subset\R^d$ będzie \emph{ograniczonym zbiorem otwartym},
a $K\Subset\Omega$ zbiorem zwartym. Oznaczmy przez $\mathcal A_\Omega(K)$ zbiór tych $v\in H_0^1(\Omega)$,
dla których istnieje otwarty $U$ spełniający $K\subset U\Subset\Omega$
oraz $v\ge1$ prawie wszędzie na $U$. Definiujemy
\begin{equation}\label{eq:pde19-cap}
 \operatorname{cap}_{2,\Omega}(K)=
 \inf_{v\in\mathcal A_\Omega(K)}
 \int_\Omega|\nabla v|^2\,\mathrm d\lambda_d(x).
\end{equation}
Pojemność zbioru pustego wynosi zero. Warunek $v\ge1$ oznacza,
że rozważamy rzeczywistą nierówność prawie wszędzie, a nie
niezdefiniowane wartości punktowe funkcji Sobolewa.
\end{definicja}

\begin{uwaga}[Pojemność i norma $H^1$]
Z nierówności Poincarégo na kostce i przedłużenia
przez zero (jak w sekcji~\ref{sec:pde18-intro}) wynika
$\|v\|_2\le C_\Omega\|\nabla v\|_2$ dla $v\in H_0^1(\Omega)$,
także bez regularności brzegu. Zatem seminorma
$\|\nabla v\|_2$ jest normą równoważną normie $H^1$
na $H_0^1(\Omega)$. Zatem zerowanie się pojemności z
\eqref{eq:pde19-cap} jest równoważne zerowaniu się pojemności
zdefiniowanej przez pełną normę $H^1$. Wartości liczbowe tych dwóch
pojemności są jednak różne. Wzory na pojemności kul poniżej dotyczą
\emph{energii gradientu}, a nie pełnej normy $H^1$.
\end{uwaga}

\begin{lemat}[Próby gładkie i obcinanie]
\label{lem:pde19-smooth}
W definicji~\ref{def:pde19-cap} można zastąpić $H_0^1(\Omega)$
przestrzenią $C_c^\infty(\Omega)$ bez zmiany wartości infimum.
Można ponadto wymagać $0\le v\le1$ prawie wszędzie, używając
wówczas funkcji Sobolewa.
\end{lemat}
\begin{proof}
Niech $v\in\mathcal A_\Omega(K)$. Obcięcie
$T(v)=\min\{1,\max\{v,0\}\}$ spełnia $T(v)=1$ prawie
wszędzie na pewnym otoczeniu $K$. Z wniosku
~\ref{cor:pde9-truncation} (stosowanego do części dodatniej
i obcięcia na poziomie $1$) wynika
$|\nabla T(v)|\le|\nabla v|$ prawie wszędzie.
Uzasadnijmy również $T(v)\in H_0^1(\Omega)$.
Dla $v_n\in C_c^\infty(\Omega)$ zbieżnych do $v$ w $H^1$
funkcje $T(v_n)$ należą do $H^1$ i mają zwarty nośnik
w $\Omega$, więc po przedłużeniu przez zero można je
wygładzić splotem o dostatecznie małym nośniku.
Należą zatem do $H_0^1(\Omega)$.
Ciąg $(T(v_n))$ jest ograniczony w $H^1$ i zbiega
do $T(v)$ w $L^2$, ponieważ $T$ jest lipschitzowska
ze stałą $1$. Słabo zbieżny podciąg w przestrzeni Hilberta
$H_0^1(\Omega)$ ma więc granicę $T(v)$.
W ten sposób obcięcie pozostaje dopuszczalne
i nie zwiększa energii.

Możemy teraz zakładać, że $0\le v\le1$ i $v=1$
prawie wszędzie na otwartym $U$ takim, że
$K\subset U\Subset\Omega$. Wybieramy
$\chi\in C_c^\infty(U)$ równą $1$ na otoczeniu $K$,
oraz $v_n\in C_c^\infty(\Omega)$ zbieżne do $v$ w $H^1$.
Kładziemy $w_n=\chi+(1-\chi)v_n$. Są to funkcje gładkie
o zwartym nośniku, równe $1$ na otoczeniu $K$.
Ponieważ $\chi(1-v)=0$ prawie wszędzie, mamy
$w_n-v=(1-\chi)(v_n-v)$ w $H^1(\Omega)$.
Reguła Leibniza i ograniczoność $\chi,\nabla\chi$
dają $w_n\to v$ w $H^1$.
W szczególności $\|\nabla w_n\|_2^2\to\|\nabla v\|_2^2$.
Infimum po funkcjach gładkich jest nie mniejsze niż
infimum po funkcjach Sobolewa, a z powyższej konstrukcji
jest nie większe od energii dowolnej dopuszczalnej funkcji
Sobolewa. Oba infima są więc równe.
\end{proof}

\begin{twierdzenie}[Monotoniczność i (skończona) podaddytywność]
\label{thm:pde19-basic}
Niech $\Omega\subset\R^d$ będzie ograniczonym zbiorem otwartym,
a $K,L\Subset\Omega$ zbiorami zwartymi. Wówczas:
\begin{enumerate}[label=\textup{(\roman*)}]
\item $0\le\operatorname{cap}_{2,\Omega}(K)<\infty$;
\item $K\subset L$ implikuje
$\operatorname{cap}_{2,\Omega}(K)\le\operatorname{cap}_{2,\Omega}(L)$;
\item $\operatorname{cap}_{2,\Omega}(K\cup L)\le
\operatorname{cap}_{2,\Omega}(K)+\operatorname{cap}_{2,\Omega}(L)$;
\item jeśli $K\Subset\Omega_1\subset\Omega_2$ i oba zbiory są
ograniczone i otwarte, to
$\operatorname{cap}_{2,\Omega_2}(K)\le
\operatorname{cap}_{2,\Omega_1}(K)$.
\end{enumerate}
Ponadto $\operatorname{cap}_{2,\Omega}(K)=0$ implikuje
$\lambda_d(K)=0$.
\end{twierdzenie}
\begin{proof}
Dla (i) wybieramy funkcję $\eta\in C_c^\infty(\Omega)$ równą
$1$ w otoczeniu $K$. Monotoniczność (ii) wynika z zawierania
zbiorów funkcji dopuszczalnych.

Niech $v,w\in H_0^1(\Omega)$ będą nieujemnymi funkcjami
dopuszczalnymi odpowiednio dla $K$ i $L$, z energiami nie większymi
od odpowiednich infimów plus $\eps$. Funkcja
$z=\max\{v,w\}=w+(v-w)^+$ należy do $H_0^1(\Omega)$:
dla części dodatniej stosujemy ten sam argument aproksymacyjny
co w lemacie~\ref{lem:pde19-smooth}. Funkcja $z$ jest
dopuszczalna dla $K\cup L$, bo jest co najmniej równa $1$
prawie wszędzie na sumie obu otoczeń dopuszczalności.
Z reguł obcięcia i zanikania gradientu na zbiorze
$\{v=w\}$ dla funkcji $v-w$ wynika, że prawie wszędzie
\[
 |\nabla\max(v,w)|^2+|\nabla\min(v,w)|^2
 =|\nabla v|^2+|\nabla w|^2,
\]
energia $z$ nie przewyższa sumy obu energii. Granica
$\eps\downarrow0$ dowodzi (iii). W (iv) przedłużamy funkcję
$H_0^1(\Omega_1)$ przez zero do $\Omega_2$, co nie zmienia
energii ani warunku dopuszczalności.

Jeżeli $v$ jest dopuszczalna dla $K$, to
$\lambda_d(K)\le\int_\Omega|v|^2\,\mathrm d\lambda_d$.
Nierówność Poincarégo (obowiązująca dla $H_0^1$ na
ograniczonym zbiorze otwartym) daje
$\lambda_d(K)\le C_\Omega\int_\Omega|\nabla v|^2$.
Biorąc infimum, otrzymujemy ostatnią tezę.
\end{proof}

\begin{uwaga}[Miara zero nie wystarcza]
Implikacji z ostatniego twierdzenia nie można odwrócić.
W wymiarze jeden pojedynczy punkt ma dodatnią pojemność,
choć oczywiście ma miarę zero. W większych wymiarach
podobnie mogą zachowywać się odpowiednio duże fragmenty
hiperpłaszczyzn (widać to z twierdzenia o śladzie).
\end{uwaga}

\section{Dokładna pojemność kondensatora kulistego}
\label{sec:pde19-ball}
\begin{twierdzenie}[Pojemność współśrodkowych kul]
\label{thm:pde19-balls}
Niech $d\ge1$, $0<r<R$ i $K=\overline{B_r(0)}$.
Dla względnej pojemności w $B_R(0)$ zachodzi
\begin{equation}\label{eq:pde19-ball-exact}
 \operatorname{cap}_{2,B_R}(K)=
 \begin{cases}
 \displaystyle\frac{2}{R-r},&d=1,\\[5pt]
 \displaystyle\frac{2\pi}{\log(R/r)},&d=2,\\[6pt]
 \displaystyle\frac{(d-2)|S^{d-1}|}{r^{2-d}-R^{2-d}},&d\ge3,
 \end{cases}
\end{equation}
gdzie $|S^{d-1}|=d\omega_d$ jest polem sfery jednostkowej.
\end{twierdzenie}
\begin{proof}
Dzięki lematowi~\ref{lem:pde19-smooth} dla dolnego
oszacowania wystarczy rozpatrzyć dopuszczalne funkcje
$v\in C_c^\infty(B_R)$. Przedłużamy je przez zero poza
kulę. Dla każdego $\vartheta\in S^{d-1}$ funkcja
$s\mapsto v(s\vartheta)$ jest gładka na $[r,R]$,
$v(r\vartheta)\ge1$ (nierówność prawie wszędzie na
otoczeniu przechodzi na punktową przez ciągłość)
i $v(R\vartheta)=0$. Z twierdzenia Newtona--Leibniza
i nierówności Cauchy'ego--Schwarza otrzymujemy
\[
 1\le\int_r^R|\partial_sv(s\vartheta)|\,\mathrm ds
 \le\left(\int_r^R|\partial_sv|^2s^{d-1}\,\mathrm ds\right)^{1/2}
 \left(\int_r^R s^{1-d}\,\mathrm ds\right)^{1/2}.
\]
Integrujemy po sferze i wykorzystujemy
$|\nabla v|\ge|\partial_sv|$. Energia całkowita jest
więc nie mniejsza od
$|S^{d-1}|/I_{r,R}$, gdzie
$I_{r,R}=\int_r^R s^{1-d}\,\mathrm ds$.

Aby uzyskać odwrotne oszacowanie, definiujemy $v=1$
na $\overline{B_r}$, $v=0$ poza $B_R$ i dla $r<s<R$
przyjmujemy funkcję radialną
\[
 v(s)=\frac{\displaystyle\int_s^R t^{1-d}\,\mathrm dt}
 {\displaystyle\int_r^R t^{1-d}\,\mathrm dt}.
\]
Jest to funkcja lipschitzowska: jej profil radialny
jest ciągły i kawałkami gładki, a pochodna
$-s^{1-d}/I_{r,R}$ na $[r,R]$ jest ograniczona.
Ma zerowy ślad na brzegu kuli, więc należy do $H_0^1(B_R)$
z twierdzenia~\ref{thm:pde10-trace-zero}.
Z rachunku radialnego otrzymujemy
\[
 \int_{B_R}|\nabla v|^2\,\mathrm d\lambda_d
 =\frac{|S^{d-1}|}{I_{r,R}^2}\int_r^R s^{1-d}\,\mathrm ds
 =\frac{|S^{d-1}|}{I_{r,R}}.
\]
Funkcja ta nie jest dopuszczalna w pierwotnej definicji,
gdyż nie jest równa co najmniej $1$ na otwartym otoczeniu
$\overline{B_r}$. Dla $0<\eps<R-r$ stosujemy jednak
ten sam wzór z $r+\eps$ zamiast $r$. Otrzymujemy
funkcję dopuszczalną o energii $|S^{d-1}|/I_{r+\eps,R}$.
Przejście $\eps\downarrow0$ daje odwrotne oszacowanie.
Co więcej, infimum w tej definicji nie jest osiągane:
otoczenie dopuszczalności każdej funkcji zawiera
$\overline{B_{r+\eps}}$ dla pewnego $\eps>0$, więc jej
energia jest co najmniej $|S^{d-1}|/I_{r+\eps,R}$,
co jest ściśle większe od $|S^{d-1}|/I_{r,R}$.
Dla $d=1$ rachunek wykonujemy na dwóch odcinkach,
a $|S^0|=2$. Mamy kolejno
$I_{r,R}=R-r$ dla $d=1$, $I_{r,R}=\log(R/r)$ dla $d=2$
oraz $I_{r,R}=(r^{2-d}-R^{2-d})/(d-2)$ dla $d\ge3$.
Daje to~\eqref{eq:pde19-ball-exact}.
\end{proof}

\begin{wniosek}[Pojemność punktu]
\label{cor:pde19-point}
Jeżeli $\Omega\subset\R^d$ jest ograniczonym zbiorem otwartym,
a $x_0\in\Omega$, to
\begin{equation}\label{eq:pde19-point}
 \operatorname{cap}_{2,\Omega}(\{x_0\})=0\quad(d\ge2),
 \qquad
 \operatorname{cap}_{2,\Omega}(\{x_0\})>0\quad(d=1).
\end{equation}
\end{wniosek}
\begin{proof}
Wybierzmy $B_R(x_0)\Subset\Omega$ i $0<r<R$.
Wzór~\eqref{eq:pde19-ball-exact} obowiązuje również dla
kul o środku $x_0$, ponieważ translacja zachowuje
dziedzinę formową, dopuszczalność i energię gradientu.
Z monotoniczności
względem zbioru i obszaru otrzymujemy
$\operatorname{cap}_{2,\Omega}(\{x_0\})\le
\operatorname{cap}_{2,B_R}(\overline{B_r(x_0)})$.
Prawa strona dąży do zera przy $r\downarrow0$ dla $d\ge2$;
w $d=2$ tempo jest logarytmiczne, a w $d\ge3$ rzędu
$r^{d-2}$. W jednym wymiarze nierówność
$|v(x_0)|\le\sqrt{|\Omega|}\,\|v'\|_{L^2(\Omega)}$
wynika następująco: dla składowej $(a,b)$ zawierającej
$x_0$ ograniczenie $v$ należy do $H_0^1(a,b)$ przez
ograniczenie ciągu aproksymujących funkcji testowych.
Jego absolutnie ciągły reprezentant spełnia $v(a)=0$,
więc $v(x_0)=\int_a^{x_0}v'(s)\,\mathrm d\lambda_1(s)$.
Stosujemy nierówność Cauchy'ego--Schwarza i $b-a\le|\Omega|$.
Dopuszczalność i ciągłość dają $v(x_0)\ge1$, skąd
$\operatorname{cap}_{2,\Omega}(\{x_0\})\ge|\Omega|^{-1}>0$.
\end{proof}

\section[Pojemność Sobolewa i wymiar krytyczny]{Pojemność $W^{1,p}$ i wymiar krytyczny}
\label{sec:pde19-p}
\begin{definicja}[Pojemność $p$-energetyczna]
Niech $1\le p<\infty$, a $K\Subset\Omega$ będzie zbiorem
zwartym, gdzie $\Omega$ jest ograniczonym zbiorem otwartym. Definiujemy
\[
 \operatorname{cap}_{p,\Omega}(K)=\inf\left\{
 \int_\Omega|\nabla v|^p\,\mathrm d\lambda_d(x):
 v\in W_0^{1,p}(\Omega),\ v\ge1\ \text{p.w. na otoczeniu }K\right\}.
\]
Otoczenie w warunku dopuszczalności można zawsze
wybrać o zwartym domknięciu w $\Omega$, dzięki zwartości
$K$. Przyjmujemy pojemność zbioru pustego równą zero.
Dla $p=2$ jest to poprzednia definicja.
\end{definicja}
\begin{twierdzenie}[Punkt a wykładnik Sobolewa]
\label{thm:pde19-point-p}
Niech $d\ge2$, $1\le p<\infty$ i $x_0$ należy do ograniczonego
zbioru otwartego $\Omega\subset\R^d$. Wtedy
$\operatorname{cap}_{p,\Omega}(\{x_0\})=0$ dla $p\le d$,
natomiast $\operatorname{cap}_{p,\Omega}(\{x_0\})>0$
dla $p>d$. W przypadku $d=1$ pojemność punktu jest dodatnia
dla każdego $p\ge1$.
\end{twierdzenie}
\begin{proof}
Po translacji przyjmijmy $x_0=0$, wybierając
$B_R\Subset\Omega$. Dla $p<d$ bierzemy
$0<r<R/2$ i próbę $v_r(x)=\eta(x/r)$, gdzie
$\eta\in C_c^\infty(B_2)$ i $\eta=1$ na $B_1$.
Mamy $v_r\in C_c^\infty(\Omega)$ i $v_r=1$ na $B_r$,
a więc funkcja jest dopuszczalna dla punktu.
Ponieważ $\nabla v_r(x)=r^{-1}\nabla\eta(x/r)$,
podstawienie $x=ry$ daje
\[
 \int_\Omega|\nabla v_r|^p\,\mathrm d\lambda_d(x)
 =r^{d-p}\int_{\R^d}|\nabla\eta(y)|^p\,\mathrm d\lambda_d(y)
 \longrightarrow0.
\]

Dla $p=d>1$ wybieramy $0<r<R$ i określamy
\[
 v_r(x)=
 \begin{cases}
 1,& |x|\le r,\\
 \log(R/|x|)/\log(R/r),&r<|x|<R,\\
 0,&|x|\ge R.
 \end{cases}
\]
Profil jest ciągły i lipschitzowski dla ustalonych $r,R$.
Funkcja ma zwarty nośnik w $\Omega$, należy do
$W^{1,d}(\R^d)$ i przez wygładzanie splotem należy
do $W_0^{1,d}(\Omega)$. Jest dopuszczalna dla punktu.
Na pierścieniu $|\nabla v_r|=1/(|x|\log(R/r))$,
poza nim gradient znika prawie wszędzie.
Jej energia $d$-ta wynosi
\[
 |S^{d-1}|\,\frac{\displaystyle\int_r^R s^{-1}\,\mathrm ds}
 {\bigl(\log(R/r)\bigr)^d}
 =\frac{|S^{d-1}|}{(\log(R/r))^{d-1}}\longrightarrow0.
\]
Dla $p>d$ przedłużamy $v\in W_0^{1,p}(\Omega)$
przez zero do $\R^d$, korzystając z twierdzenia
~\ref{thm:pde9-zero-extension}. Wybieramy ustalony
punkt $y\notin\overline\Omega$. Nierówność Morreya
~\ref{thm:pde10-Morrey} dla ciągłego reprezentanta
przedłużenia $E_0v$ daje
\[
 |v(x_0)|=|E_0v(x_0)-E_0v(y)|
 \le C_{d,p}|x_0-y|^{1-d/p}\|\nabla v\|_{L^p(\Omega)}.
\]
Równość $E_0v(y)=0$ wynika z ciągłości i zerowania się
prawie wszędzie na otoczeniu $y$.
Dopuszczalność oraz ciągłość dają $v(x_0)\ge1$.
Podnosząc oszacowanie do potęgi $p$, otrzymujemy
dolną granicę energii niezależną od $v$, ściśle dodatnią.

W $d=1$ dla każdego $1\le p<\infty$ stosujemy
absolutną ciągłość na składowej $(a,b)$, jak w dowodzie
wniosku~\ref{cor:pde19-point}. Dla $p>1$ nierówność
Höldera daje
$1\le |v(x_0)|\le (b-a)^{1-1/p}\|v'\|_{L^p(a,b)}$.
Dla $p=1$ stosujemy bezpośrednio
$1\le\int_a^b|v'|\,\mathrm d\lambda_1$.
W obu przypadkach energia jest co najmniej
$(b-a)^{1-p}>0$.
\end{proof}

\begin{uwaga}[Wymiar, miara i pojemność]
Wynik nie orzeka, że każdy zbiór wymiaru Hausdorffa $d-p$
ma zerową pojemność: przypadki graniczne są subtelne.
Zależności z miarą Hausdorffa wymagają osobnych twierdzeń,
a dla $p=d$ pojawiają się poprawki logarytmiczne.
\end{uwaga}

\section{Usuwalność osobliwości dla słabych rozwiązań}
\label{sec:pde19-removable}
\begin{twierdzenie}[Zbiór zerowej pojemności jest niewidoczny dla rozwiązań energetycznych]
\label{thm:pde19-remove}
Niech $\Omega\subset\R^d$ będzie ograniczonym zbiorem otwartym,
a $K\Subset\Omega$ zbiorem zwartym takim, że
$\operatorname{cap}_{2,\Omega}(K)=0$.
Niech $A\in L^\infty(\Omega;\R^{d\times d})$;
nie zakładamy tu ani symetrii, ani eliptyczności.
Jeżeli $u\in H^1(\Omega\setminus K)$ oraz
\begin{equation}\label{eq:pde19-off-K}
 \int_{\Omega\setminus K}A(x)\nabla u\cdot\nabla\varphi
 \,\mathrm d\lambda_d(x)=0
 \quad\text{dla każdego }\varphi\in C_c^\infty(\Omega\setminus K),
\end{equation}
to $u$ ma jedyne (z dokładnością do równości prawie
wszędzie) przedłużenie do $H^1(\Omega)$, a przedłużenie
jest słabym rozwiązaniem $-\operatorname{div}(A\nabla u)=0$
na całym $\Omega$.
\end{twierdzenie}
\begin{proof}
Z lematu~\ref{lem:pde19-smooth} i zerowej pojemności
istnieją $\eta_j\in C_c^\infty(\Omega)$, $\eta_j\ge1$
na otoczeniu $K$, takie że $\|\nabla\eta_j\|_2\to0$.
Wybieramy gładką funkcję $\Theta\colon\R\to[0,1]$,
równą $0$ na $(-\infty,0]$ i $1$ na $[1,\infty)$,
z ograniczoną pochodną (np. całkę z nieujemnej
funkcji testowej na $(0,1)$ o całce $1$). Wtedy
$\psi_j=\Theta(\eta_j)\in C_c^\infty(\Omega)$,
$\psi_j=1$ w otoczeniu $K$ i
$\|\psi_j\|_{H^1(\Omega)}\to0$: ostatni wniosek wynika
z nierówności Poincarégo i oszacowań
$|\Theta(s)|\le C|s|$, $|\Theta'(s)|\le C$.

Z twierdzenia~\ref{thm:pde19-basic} mamy $\lambda_d(K)=0$.
Przedłużamy $u$ oraz każdą jego słabą pochodną
$g_\ell=\partial_\ell u$ przez zero na $K$ jako funkcje
$L^2(\Omega)$. Aby wykazać, że przedłużone $g_\ell$
są słabymi pochodnymi przedłużonego $u$, dla
$\varphi\in C_c^\infty(\Omega)$ stosujemy równanie
definiujące słabą pochodną na $\Omega\setminus K$
do testu $(1-\psi_j)\varphi$. Po rozwinięciu otrzymujemy
\[
 \int_\Omega u\,\partial_\ell\varphi\,\mathrm d\lambda_d
 +\int_\Omega g_\ell\varphi\,\mathrm d\lambda_d
 =
 \int_\Omega (u\psi_j\partial_\ell\varphi+
 u\varphi\partial_\ell\psi_j+g_\ell\psi_j\varphi)
 \,\mathrm d\lambda_d.
\]
Moduł prawej strony jest nie większy od
\[
 (\|u\|_2\|\partial_\ell\varphi\|_\infty+
 \|g_\ell\|_2\|\varphi\|_\infty)\|\psi_j\|_2+
 \|u\|_2\|\varphi\|_\infty\|\partial_\ell\psi_j\|_2,
\]
co dąży do zera. Zatem $u\in H^1(\Omega)$ po
przedłużeniu i jego gradient zgadza się prawie wszędzie
z gradientem na $\Omega\setminus K$.
Jednoznaczność wynika z $\lambda_d(K)=0$.

Weźmy teraz dowolną $\varphi\in C_c^\infty(\Omega)$.
Funkcja $\varphi_j=(1-\psi_j)\varphi$ jest gładka,
ma zwarty nośnik w $\Omega\setminus K$, więc z
\eqref{eq:pde19-off-K} mamy
\[
 0=\int_\Omega A\nabla u\cdot\nabla\varphi_j
 \mathrm d\lambda_d
 =\int_\Omega A\nabla u\cdot\nabla\varphi
 \mathrm d\lambda_d
 -\int_\Omega A\nabla u\cdot
 (\psi_j\nabla\varphi+\varphi\nabla\psi_j)\mathrm d\lambda_d.
\]
Dwa ostatnie składniki dążą do zera dzięki
nierówności Cauchy'ego--Schwarza, ograniczoności
$A,\varphi,\nabla\varphi$ i zbieżności
$\psi_j\to0$ w $H^1$. Otrzymujemy tożsamość słabą
na całym $\Omega$.
\end{proof}

\begin{wniosek}[Równość przestrzeni Dirichleta po usunięciu zbioru]
\label{cor:pde19-H0}
Niech $\Omega$ będzie ograniczonym zbiorem otwartym,
a $K\Subset\Omega$ zbiorem zwartym o
$\operatorname{cap}_{2,\Omega}(K)=0$. Wtedy zachodzi
\[
 H_0^1(\Omega\setminus K)=H_0^1(\Omega)
\]
po naturalnym utożsamieniu funkcji przez przedłużenie zerem
na $K$ (który ma miarę zero).
\end{wniosek}
\begin{proof}
Z twierdzenia~\ref{thm:pde19-basic} zbiór $K$ ma miarę
Lebesgue'a zero. Przedłużenie przez zero funkcji
testowych z $\Omega\setminus K$ do $\Omega$ zachowuje
normę $H^1$. Przez gęstość daje ono izometryczne
zanurzenie $H_0^1(\Omega\setminus K)\to H_0^1(\Omega)$,
którego obraz jest domknięty dzięki zupełności.
Dla odwrotnego zawierania wystarczy więc wykazać,
że obraz zawiera $C_c^\infty(\Omega)$.
Jeżeli $\varphi\in C_c^\infty(\Omega)$ i $\psi_j$ są jak
w poprzednim dowodzie, to
$\varphi_j=(1-\psi_j)\varphi\in C_c^\infty(\Omega\setminus K)$
i $\varphi_j\to\varphi$ w $H^1(\Omega)$, ponieważ
$\|\varphi\psi_j\|_2\le\|\varphi\|_\infty\|\psi_j\|_2$
oraz
$\|\nabla(\varphi\psi_j)\|_2\le
\|\nabla\varphi\|_\infty\|\psi_j\|_2+
\|\varphi\|_\infty\|\nabla\psi_j\|_2$.
Z domkniętości obrazu i gęstości $C_c^\infty(\Omega)$
wynika żądana równość przestrzeni.
\end{proof}

\begin{kontrprzyklad}[Zerowa pojemność nie wystarcza bez kontroli energii]
\label{ex:pde19-fundamental}
Dla $d\ge3$ funkcja $u(x)=|x|^{2-d}$, a dla $d=2$
$u(x)=\log|x|$, jest harmoniczna na $B_1\setminus\{0\}$,
ale nie należy do $H^1(B_1\setminus\{0\})$ ani nie
ma przedłużenia do $H^1(B_1)$. Rzeczywiście,
$|\nabla u(x)|\asymp|x|^{1-d}$ (dla $d\ge2$), więc
\[
 \int_{B_\eps\setminus\{0\}}|\nabla u|^2
 \,\mathrm d\lambda_d(x)
 \asymp\int_0^\eps r^{1-d}\,\mathrm dr=+\infty.
\]
Funkcja $u$ jest lokalnie całkowalna, gdyż dla $d\ge3$
całka radialna modułu w pobliżu zera sprowadza się
do $\int_0^\eps r\,\mathrm dr$, a dla $d=2$
do $\int_0^\eps r|\log r|\,\mathrm dr$.
W sensie dystrybucyjnym, przy normalizacji rozwiązania
fundamentalnego z~\eqref{eq:pde7-laplace-fundamental}, mamy
\[
 -\Delta u=
 \begin{cases}
 (d-2)|S^{d-1}|\delta_0,&d\ge3,\\
 -2\pi\delta_0,&d=2.
 \end{cases}
\]
W szczególności przedłużenie lokalnie całkowalne
nie jest harmoniczne w sensie dystrybucji. Zerowa pojemność punktu nie
usuwa więc osobliwości przy nieograniczonej energii.
\end{kontrprzyklad}

\section{Funkcje harmoniczne ograniczone na obszarze z przekłuciem}
\label{sec:pde19-puncture}
\begin{twierdzenie}[Usuwalność ograniczonej osobliwości izolowanej]
\label{thm:pde19-bounded}
Niech $d\ge2$, $B_R(x_0)\subset\R^d$ będzie kulą,
a $u\in C^2(B_R(x_0)\setminus\{x_0\})$ funkcją harmoniczną.
Jeżeli $u$ jest ograniczona w otoczeniu $x_0$, to istnieje
jedyna funkcja harmoniczna $\widetilde u\in C^\infty(B_R(x_0))$
równa $u$ poza $x_0$.
\end{twierdzenie}
\begin{proof}
Po translacji przyjmijmy $x_0=0$. Wybierzmy
$\chi\in C_c^\infty(B_R)$, $0\le\chi\le1$, równą $1$
na $B_{R/2}$. Ustalmy promienie $R_j\downarrow0$,
$R_j<R/2$. Z wniosku~\ref{cor:pde19-point} pojemność
punktu w każdej kuli $B_{R_j}$ wynosi zero. Z lematu
~\ref{lem:pde19-smooth} oraz gładkiego obcięcia
z dowodu twierdzenia~\ref{thm:pde19-remove} wybieramy
$\psi_j\in C_c^\infty(B_{R_j})$ równe $1$ na
otoczeniu zera, $0\le\psi_j\le1$, tak aby
$\|\nabla\psi_j\|_{L^2(\R^d)}\le1/j$.
Po przedłużeniu przez zero mamy $\psi_j(x)\to0$
dla każdego $x\ne0$.

Funkcja $u$ jest ograniczona na
$\operatorname{supp}\chi\setminus\{0\}$:
w pobliżu zera z założenia, a poza mniejszą kulą
z ciągłości na zbiorze zwartym. Niech $M_0$ będzie
kresem $|u|$ na tym zbiorze. Kładziemy
$\zeta_j=\chi(1-\psi_j)$. Funkcja $\zeta_j^2u$
ma zwarty nośnik poza zerem i należy do $H^1$;
przez lokalne wygładzanie jest dopuszczalnym testem
w równaniu harmonicznym. Otrzymujemy
\[
 \int\zeta_j^2|\nabla u|^2\,\mathrm d\lambda_d
 =-2\int \zeta_j u\,\nabla u\cdot\nabla\zeta_j\,
 \mathrm d\lambda_d.
\]
Całki liczymy na $B_R\setminus\{0\}$.
Z nierówności Cauchy'ego--Schwarza wynika
$\int\zeta_j^2|\nabla u|^2\,\mathrm d\lambda_d
\le4\int |u|^2|\nabla\zeta_j|^2\,\mathrm d\lambda_d$
(po podzieleniu przez pierwiastek z lewej strony,
gdy jest dodatni; przypadek zerowy jest natychmiastowy).
Ponieważ $\nabla\zeta_j=(1-\psi_j)\nabla\chi-\chi\nabla\psi_j$,
otrzymujemy
\[
 \int_{B_R}\chi^2(1-\psi_j)^2|\nabla u|^2
 \,\mathrm d\lambda_d
 \le 8 M_0^2\int_{B_R}
 (|\nabla\chi|^2+\chi^2|\nabla\psi_j|^2)
 \,\mathrm d\lambda_d,
\]
Prawa strona jest jednostajnie ograniczona względem $j$.
Lemat Fatou i $\psi_j\to0$ poza zerem dają
$\nabla u\in L^2(B_{R/2}\setminus\{0\})$.
Ponadto $u\in L^2(B_{R/2}\setminus\{0\})$ z ograniczoności.
Ponieważ $u$ jest klasycznie gładka poza zerem (wystarcza
tu $C^2$), jej klasyczne pochodne są tam słabymi
pochodnymi. Zatem $u\in H^1(B_{R/2}\setminus\{0\})$.
Twierdzenie~\ref{thm:pde19-remove}, zastosowane z $A=I$,
daje przedłużenie do $H^1(B_{R/2})$ harmoniczne
dystrybucyjnie. Lemat Weyla~\ref{thm:pde7-weyl}
daje jego gładkiego harmonicznego reprezentanta.
Na przekłutej kuli $B_{R/2}\setminus\{0\}$ zgadza się on
z pierwotną funkcją wszędzie, bo obie są ciągłe i równe
prawie wszędzie. Możemy więc skleić go z $u$ na
$B_R\setminus\{0\}$, otrzymując funkcję harmoniczną
na $B_R$. Jednoznaczność wynika z ciągłości
i gęstości $B_R\setminus\{0\}$ w $B_R$.
\end{proof}

\begin{przyklad}[Punkt brzegowy, w którym dane Dirichleta giną]
\label{ex:pde19-irregular}
Niech $d\ge2$ i $\Omega=B_1(0)\setminus\{0\}$.
Na brzegu $\partial\Omega=S^{d-1}\cup\{0\}$ zadajmy
funkcję ciągłą $g$, równą $0$ na $S^{d-1}$ oraz $1$ w zerze.
Nie istnieje ograniczona funkcja harmoniczna na $\Omega$,
ciągła na $\overline\Omega$ i przyjmująca wartości $g$.
W istocie taka funkcja na mocy poprzedniego twierdzenia
przedłużyłaby się harmonicznie na $B_1$. Z zasady maksimum~\ref{thm:pde14-weak-max},
stosowanej do przedłużenia oraz do jego przeciwności
przy zerowych danych na $S^{d-1}$, byłaby wszędzie zerowa.
Przeczy to wartości $1$ narzuconej w zerze.
To przykład \emph{nieregularnego punktu brzegu}.
\end{przyklad}

\section{Pojemność, reprezentanci i zbiory wyjątkowe}
\label{sec:pde19-quasi}
\begin{definicja}[Pojemność zbioru otwartego i quasi-ciągłość]
Niech $\Omega$ będzie ograniczonym zbiorem otwartym.
Dla otwartego $U\subset\Omega$ definiujemy
\[
 \operatorname{cap}_{2,\Omega}(U)=\inf\left\{
 \|\nabla v\|_2^2:v\in H_0^1(\Omega),\
 v\ge1\ \text{p.w. na }U\right\}.
\]
Nie wymagamy $U\Subset\Omega$ i przyjmujemy
$\inf\varnothing=+\infty$.
Dla dowolnego $E\subset\Omega$ kładziemy
\[
 \operatorname{cap}_{2,\Omega}(E)
 =\inf\{\operatorname{cap}_{2,\Omega}(U):E\subset U\subset\Omega,
 \ U\text{ otwarty}\}.
\]
Rozszerzenie jest zgodne z definicją dla zbiorów
otwartych dzięki monotoniczności warunku dopuszczalności,
a dla zbiorów zwartych $K\Subset\Omega$ dzięki możliwości
zmniejszenia otwartego otoczenia do otoczenia
o zwartym domknięciu w $\Omega$.
Funkcję $\widetilde u:\Omega\to\R$ nazywamy
\emph{quasi-ciągłą}, jeśli dla każdego $\eps>0$
istnieje otwarty $G\subset\Omega$,
$\operatorname{cap}_{2,\Omega}(G)<\eps$, taki że
$\widetilde u|_{\Omega\setminus G}$ jest ciągła.
Własność zachodzącą poza zbiorem pojemności zero nazywamy
własnością zachodzącą \emph{quasi-wszędzie}.
\end{definicja}

\begin{uwaga}
Rozszerzona pojemność jest monotoniczna. Dla otwartego
$U$ nierówność Poincarégo daje
$\lambda_d(U)\le C_\Omega^2\operatorname{cap}_{2,\Omega}(U)$
(przez oszacowanie każdej próby, również gdy infimum
wynosi $+\infty$). Dla dowolnego $E\subset\Omega$ stąd
$\lambda_d^*(E)\le C_\Omega^2\operatorname{cap}_{2,\Omega}(E)$,
gdzie $\lambda_d^*$ jest miarą zewnętrzną Lebesgue'a.
W szczególności zbiory pojemności zero mają miarę zero.
\end{uwaga}

\begin{lemat}[Podaddytywność przeliczalna pojemności]
\label{lem:pde19-countable}
Dla dowolnych $E_j\subset\Omega$, $j\ge1$, zachodzi
$\operatorname{cap}_{2,\Omega}(\bigcup_jE_j)
 \le\sum_j\operatorname{cap}_{2,\Omega}(E_j)$.
\end{lemat}
\begin{proof}
Gdy suma pojemności po prawej stronie nierówności
jest nieskończona, teza jest natychmiastowa.
Najpierw rozważmy zbiory otwarte $U_j$ o skończonej sumie
pojemności i ustalmy $\eps>0$. Wybieramy $v_j\in H_0^1(\Omega)$, $v_j\ge1$
prawie wszędzie na $U_j$, $0\le v_j\le1$, z
$\int_\Omega|\nabla v_j|^2\,\mathrm d\lambda_d(x)\le\operatorname{cap}_{2,\Omega}(U_j)
 +\eps 2^{-j}$. Dla $m\ge1$ kładziemy
$w_m=\max(v_1,\ldots,v_m)$. Z tożsamości gradientowej
użytej w dowodzie Twierdzenia~\ref{thm:pde19-basic}
wynika $\|\nabla w_m\|_2^2\le\sum_{j=1}^m
\|\nabla v_j\|_2^2$. Normy $H_0^1$ są więc jednostajnie
ograniczone. Ponadto $0\le w_m\le1$ i $w_m$ jest
niemalejący prawie wszędzie. Po odrzuceniu przeliczalnej
sumy zbiorów miary zero wybieramy punktową granicę $w$.
Z twierdzenia o zbieżności zmajoryzowanej zastosowanego
do $|w_m-w|^2\le1$ na ograniczonym zbiorze $\Omega$
wynika $w_m\to w$ silnie w $L^2$.
Z ograniczoności w przestrzeni Hilberta $H_0^1$
wybieramy słabo zbieżny podciąg; jego granica musi być
równa $w$ przez ciągłe zanurzenie do $L^2$.
Słaba dolna półciągłość normy gradientu daje
$w\in H_0^1$ oraz
$\|\nabla w\|_2^2\le\sum_j\|\nabla v_j\|_2^2$.
Ponieważ $w\ge1$ prawie wszędzie na $\bigcup_jU_j$,
uzyskujemy
$\operatorname{cap}_{2,\Omega}(\bigcup_jU_j)
\le\sum_j\operatorname{cap}_{2,\Omega}(U_j)+\eps$.
Przejście $\eps\downarrow0$ dowodzi nierówności
dla zbiorów otwartych.
Z definicji zewnętrznej pojemności wybieramy
$U_j\supset E_j$ z
$\operatorname{cap}_{2,\Omega}(U_j)
\le\operatorname{cap}_{2,\Omega}(E_j)+\eps2^{-j}$.
Stosujemy monotoniczność i już udowodnioną nierówność
dla $\bigcup_jU_j$, otrzymując górną granicę
$\sum_j\operatorname{cap}_{2,\Omega}(E_j)+\eps$.
Przechodzimy z $\eps\downarrow0$.
\end{proof}

\begin{twierdzenie}[Quasi-ciągły reprezentant funkcji Sobolewa]
\label{thm:pde19-quasi}
Niech $\Omega\subset\R^d$ będzie ograniczonym zbiorem otwartym.
Każda rzeczywista funkcja $u\in H_0^1(\Omega)$ ma
reprezentanta quasi-ciągłego $\widetilde u$, równego $u$
prawie wszędzie. Jest on jednoznaczny z dokładnością
do zbioru pojemności zero.
\end{twierdzenie}
\begin{proof}
Wybieramy $u_n\in C_c^\infty(\Omega)$ tak szybko zbieżne
do $u$ w $H_0^1(\Omega)$, aby
$\|u_n-u\|_{H^1}\le2^{-2n-1}$ dla $n\ge1$.
Z nierówności trójkąta wynika wtedy
$\|u_{n+1}-u_n\|_{H^1}^2\le2^{-4n}$.
Niech $G_n=\{x\in\Omega:|u_{n+1}(x)-u_n(x)|>2^{-n}\}$.
Ponieważ $u_{n+1}-u_n\in C_c^\infty(\Omega)$,
funkcja $2^n|u_{n+1}-u_n|$ należy do $H_0^1$ i jest
co najmniej $1$ na $G_n$. Zatem
\[
 \operatorname{cap}_{2,\Omega}(G_n)
 \le2^{2n}\|\nabla(u_{n+1}-u_n)\|_2^2\le2^{-2n}.
\]
Dla $F_m=\bigcup_{n\ge m}G_n$ z
lematu~\ref{lem:pde19-countable} mamy
$\operatorname{cap}_{2,\Omega}(F_m)
\le\sum_{n\ge m}2^{-2n}\to0$.
Na $\Omega\setminus F_m$ szereg
$\sum_{n\ge m}(u_{n+1}-u_n)$ jest jednostajnie i bezwzględnie
zbieżny, ponieważ każdy składnik ma tam moduł nie większy niż
$2^{-n}$. Granica $\widetilde u$ jest na tym zbiorze ciągła
jako jednostajna granica funkcji ciągłych.

Zbiór $F=\bigcap_m F_m$ ma pojemność zero, ponieważ
$F\subset F_m$ dla wszystkich $m$; poza $F$
ciąg $u_n$ zbiega punktowo. Ponieważ $F$ ma miarę zero,
a ze zbieżności w $L^2$ można wybrać podciąg zbieżny
prawie wszędzie do $u$, otrzymujemy $\widetilde u=u$
prawie wszędzie. Wartości na $F$ można zdefiniować
dowolnie. Dla każdego $\eps>0$ wybór dużego $m$ sprawia,
że $\operatorname{cap}_{2,\Omega}(F_m)<\eps$,
a ograniczenie $\widetilde u$ do dopełnienia $F_m$
jest ciągłe.

Dla jednoznaczności niech $v,w$ będą dwoma
quasi-ciągłymi reprezentantami tej samej klasy.
Dla $\eps>0$ wybieramy otwarte zbiory wyjątkowe
dla obu funkcji, każdy o pojemności mniejszej niż
$\eps/2$. Ich suma $G$ ma pojemność mniejszą niż
$\eps$ i obie funkcje są ciągłe na $\Omega\setminus G$.
Niech $D=\{x:v(x)\ne w(x)\}$. Jest to zbiór zawarty
w zbiorze miary zero. Z względnej ciągłości wynika,
że $D\setminus G$ jest otwarty w $\Omega\setminus G$,
czyli istnieje otwarty $V\subset\Omega$ taki, że
$D\setminus G=V\setminus G$. Zatem $D\subset V\cup G$,
a $\lambda_d(V\setminus G)=0$. Każda próba dopuszczalna
dla $G$ jest dopuszczalna także dla otwartego $V\cup G$,
gdyż warunek obowiązuje prawie wszędzie.
Stąd $\operatorname{cap}_{2,\Omega}(V\cup G)
=\operatorname{cap}_{2,\Omega}(G)<\eps$.
Monotoniczność daje $\operatorname{cap}_{2,\Omega}(D)<\eps$,
a dowolność $\eps$ dowodzi jednoznaczności.
\end{proof}

\begin{uwaga}[Po co quasi-ciągłość?]
Funkcję Sobolewa określamy zwykle jedynie prawie wszędzie,
więc zapis $u(x_0)$ nie ma sensu dla dowolnie wybranego
punktu. Twierdzenie~\ref{thm:pde19-quasi} zastępuje
nieistniejącą ciągłość klasyczną ciągłością poza zbiorem
\emph{małym energetycznie}. Jest to podstawa subtelnej
interpretacji warunków brzegowych oraz przeszkód w zagadnieniach
wariacyjnych. Nie utożsamiamy przy tym automatycznie
pojęcia quasi-ciągłości z regularnym śladem na
brzegu Lipschitzowskim.
\end{uwaga}

\section{Bariera zewnętrznej kuli i kryterium Wienera}
\label{sec:pde19-wiener}
\begin{definicja}[Bariera harmoniczna]
Niech $\Omega\subset\R^d$ będzie ograniczonym obszarem,
a $x_0\in\partial\Omega$. Lokalną barierą harmoniczną
w $x_0$ nazywamy nieujemną funkcję harmoniczną $b$
w $\Omega\cap B_\rho(x_0)$, dodatnią w jego wnętrzu,
która spełnia $b(x)\to0$ przy $x\to x_0$ od strony $\Omega$.
Do kryterium regularności używamy bariery spełniającej
dodatkowo warunek oddzielenia od zera poza każdym
otoczeniem $x_0$: dla $0<\delta<\rho$ infimum $b$
na $\Omega\cap B_\rho(x_0)\setminus B_\delta(x_0)$
jest dodatnie, jeśli ten zbiór jest niepusty.
W szczególności nie wystarcza sama dodatniość
we wnętrzu.
\end{definicja}

\begin{stwierdzenie}[Jawna bariera przy zewnętrznej kuli]
\label{prop:pde19-exterior-ball}
Niech $d\ge2$, $\Omega\subset\R^d$ będzie ograniczonym
obszarem, a $x_0\in\partial\Omega$. Załóżmy, że istnieje
kula $B_r(y)\subset\R^d\setminus\overline\Omega$
taka, że $|x_0-y|=r$. Wtedy funkcja
\[
 b(x)=\begin{cases}
 r^{2-d}-|x-y|^{2-d},&d\ge3,\\[2pt]
 \log(|x-y|/r),&d=2
 \end{cases}
\]
jest dodatnia i harmoniczna w $\Omega$, a
$b(x)\to0$ przy $x\to x_0$ w $\Omega$.
\end{stwierdzenie}
\begin{proof}
Ponieważ kula $B_r(y)$ jest zawarta w dopełnieniu
$\overline\Omega$, każdy $x\in\Omega$ spełnia
$|x-y|>r$. Funkcja radialna $s\mapsto s^{2-d}$
(dla $d\ge3$) albo $s\mapsto\log s$ (dla $d=2$)
jest harmoniczna poza środkiem $y$. Wskazane kombinacje
są dodatnie dla $s>r$ i zerują się dla $s=r$.
Ciągłość funkcji radialnej daje granicę zero w $x_0$.
Jeśli kula styka się z brzegiem także w innych punktach,
sama ta funkcja może nie spełniać warunku oddzielenia
od zera. Wybierzmy $0<r'<r$ oraz
$y'=x_0+(r'/r)(y-x_0)$. Z nierówności trójkąta
$\overline{B_{r'}(y')}\subset\overline{B_r(y)}$.
Równość w nierówności trójkąta pokazuje, że sfery
obu kul mają tylko punkt $x_0$ wspólny: dla
$|x-y'|\le r'$ i $|x-y|=r$ konieczne jest
$x-y'$ współliniowe z $y'-y$ i zgodnie skierowane,
o długości $r'$, co daje $x=x_0$.
W konsekwencji mniejsza domknięta kula przecina
$\overline\Omega$ wyłącznie w $x_0$. Wstawiając
$r',y'$ do wzoru na $b$, otrzymujemy funkcję ciągłą
na $\overline\Omega$, dodatnią na
$\overline\Omega\setminus\{x_0\}$ i zerową w $x_0$.
Dla każdego $\delta>0$ jej minimum na niepustym
zwartym zbiorze $\overline\Omega\setminus B_\delta(x_0)$
jest ściśle dodatnie. Jest to bariera z wymaganym
warunkiem oddzielenia.
\end{proof}

\begin{uwaga}[Warunek geometryczny a rozwiązywalność]
Istnienie bariery zewnętrznej kuli prowadzi do regularności
$x_0$ dla zagadnienia Dirichleta w sensie Perrona, po
zastosowaniu zasady porównania. Korzystamy tu z klasycznego twierdzenia o barierach
dla rozwiązania Perrona (wynik przywołany, nieudowodniony
w tym rozdziale). Pełna konstrukcja Perrona
nie jest tutaj dowodzona. Warunek zewnętrznej kuli jest
wystarczający, ale niekonieczny.
\end{uwaga}

\begin{definicja}[Bezwzględna pojemność Newtonowska]
Dla $d\ge3$ i zwartego $K\subset\R^d$ określamy
\[
 \operatorname{Cap}_2(K)=\inf\left\{
 \int_{\R^d}|\nabla\varphi|^2\,\mathrm d\lambda_d(x):
 \varphi\in C_c^\infty(\R^d),\ \varphi\ge1
 \text{ na otoczeniu }K\right\}.
\]
W odróżnieniu od~\eqref{eq:pde19-cap} nie występuje tu
ustalony obszar zewnętrzny ani wyraz $L^2$.
\end{definicja}

\begin{stwierdzenie}[Skalowanie i pojemność bezwzględna kuli]
Dla $d\ge3$, zwartego $K\subset\R^d$, $a>0$ i
$z\in\R^d$ zachodzi
\[
 \operatorname{Cap}_2(z+aK)=a^{d-2}\operatorname{Cap}_2(K),
 \qquad
 \operatorname{Cap}_2(\overline{B_\rho(z)})
 =(d-2)|S^{d-1}|\rho^{d-2}\quad(\rho>0).
\]
Pojemność jest monotoniczna względem zawierania zbiorów.
\end{stwierdzenie}
\begin{proof}
Monotoniczność wynika z odwrotnego zawierania klas
dopuszczalnych. Przyporządkowanie
$\varphi\mapsto\varphi((\cdot-z)/a)$ jest bijekcją
między klasami prób dla $K$ i $z+aK$. Zmiana
zmiennych w energii gradientu daje czynnik $a^{d-2}$.
Dla kuli o środku zero każda próba gładka o zwartym
nośniku należy do $C_c^\infty(B_R)$ dla pewnego
$R>\rho$. Odwrotnie, każda taka próba jest również
próbą w całej przestrzeni po przedłużeniu przez zero.
Z lematu~\ref{lem:pde19-smooth} wynika zatem
\[
 \operatorname{Cap}_2(\overline{B_\rho})
 =\inf_{R>\rho}\operatorname{cap}_{2,B_R}(\overline{B_\rho})
 =\lim_{R\to\infty}
 \frac{(d-2)|S^{d-1}|}{\rho^{2-d}-R^{2-d}}.
\]
Ostatnią równość daje monotoniczność względem obszaru
i wzór~\eqref{eq:pde19-ball-exact}.
Granica wynosi $(d-2)|S^{d-1}|\rho^{d-2}$.
Przesunięcie środka nie zmienia pojemności.
\end{proof}


\begin{lemat}[Porównanie pojemności względnej i bezwzględnej]
\label{lem:pde19-cap-comparison}
Niech $d\ge3$, $\Omega\subset\R^d$ będzie ograniczonym
zbiorem otwartym, a $D\Subset\Omega$ zwartym zbiorem.
Istnieje stała $C=C(d,D,\Omega)$ taka, że dla każdego
zwartego $K\subset D$
\[
 \operatorname{Cap}_2(K)\le\operatorname{cap}_{2,\Omega}(K)
 \le C\operatorname{Cap}_2(K).
\]
W szczególności zerowość tych dwóch pojemności jest
równoważna na ustalonym zwartym podzbiorze wnętrza $\Omega$.
\end{lemat}
\begin{proof}
W pierwszej nierówności stosujemy przedłużenie przez
zero funkcji dopuszczalnych. Dzięki
Lematowi~\ref{lem:pde19-smooth} możemy ograniczyć się
do prób gładkich. Dla drugiej nierówności wybierzmy
$\chi\in C_c^\infty(\Omega)$, $\chi=1$ na otoczeniu $D$.
Jeżeli $\varphi\in C_c^\infty(\R^d)$ jest dopuszczalna
dla $\operatorname{Cap}_2(K)$, to $\chi\varphi$ jest
dopuszczalna dla $\operatorname{cap}_{2,\Omega}(K)$.
Mamy $|\nabla(\chi\varphi)|^2\le
2\chi^2|\nabla\varphi|^2+2|\nabla\chi|^2|\varphi|^2$.
Niech $S=\operatorname{supp}\nabla\chi$ oraz
$2^*=2d/(d-2)$. Z nierówności Höldera
i twierdzenia Sobolewa~\ref{thm:pde10-Sobp} dla $p=2$
otrzymujemy
\[
 \int_S|\varphi|^2\,\mathrm d\lambda_d
 \le |S|^{2/d}\|\varphi\|_{L^{2^*}(\R^d)}^2
 \le C_d^2|S|^{2/d}\|\nabla\varphi\|_{L^2(\R^d)}^2.
\]
Po uwzględnieniu ograniczoności $\chi$ i $\nabla\chi$
wynika
\[
 \int_\Omega|\nabla(\chi\varphi)|^2
 \,\mathrm d\lambda_d(x)
 \le C\int_{\R^d}|\nabla\varphi|^2
 \,\mathrm d\lambda_d(x).
\]
Bierzemy infimum po $\varphi$.
\end{proof}

\begin{twierdzenie}[Kryterium Wienera, $d\ge3$ — wynik przywołany]
\label{thm:pde19-wiener}
Niech $d\ge3$, $\Omega\subset\R^d$ będzie ograniczonym
obszarem i $x_0\in\partial\Omega$. Dla $k$ dostatecznie
dużych niech
\[
 E_k=(\R^d\setminus\Omega)\cap
 \{x:2^{-k-1}\le|x-x_0|\le2^{-k}\}.
\]
Punkt $x_0$ jest regularny dla klasycznego zagadnienia
Dirichleta w sensie Perrona wtedy i tylko wtedy, gdy
\begin{equation}\label{eq:pde19-wiener}
 \sum_{k=k_0}^\infty
 2^{k(d-2)}\operatorname{Cap}_2(E_k)=+\infty.
\end{equation}
Tutaj regularność oznacza, że rozwiązanie Perrona
przyjmuje w granicy od strony $\Omega$ wartość
$f(x_0)$ dla każdej ciągłej funkcji brzegowej
$f\in C(\partial\Omega)$.
\end{twierdzenie}

\begin{uwaga}[Zakres dowodu kryterium]
Kryterium Wienera jest głębokim twierdzeniem teorii
potencjału: \emph{nie dowodzimy tu jego pełnej implikacji
w obu kierunkach}. Z poprzednich twierdzeń wykazaliśmy
natomiast jawnie wzór pojemności kul i nieregularność
izolowanego punktu brzegowego. Dla przekłutej kuli
z~\ref{ex:pde19-irregular} wszystkie dostatecznie małe
$E_k$ są puste, zatem szereg jest zbieżny. Dla dopełnienia
zawierającego stały stożek z wierzchołkiem w $x_0$
pojemności kolejnych fragmentów są porównywalne z
$2^{-k(d-2)}$, co daje rozbieżność. Kryterium w wymiarze
dwa przyjmuje postać logarytmiczną i nie jest tu
formułowane.
\end{uwaga}


\begin{stwierdzenie}[Zewnętrzny stożek a szereg Wienera]
\label{prop:pde19-cone-wiener}
Niech $d\ge3$, $\Omega\subset\R^d$ będzie ograniczonym obszarem
oraz $x_0\in\partial\Omega$. Załóżmy, że istnieją $r_0>0$
i niepusty, otwarty zbiór kierunków $\Sigma\subset S^{d-1}$
takie, że $x_0+r\theta\notin\Omega$ dla każdego
$0<r<r_0$ i $\theta\in\Sigma$.
Wtedy dla zbiorów $E_k$ z Twierdzenia~\ref{thm:pde19-wiener}
istnieje $c>0$ i $k_1$ takie, że
\[
  2^{k(d-2)}\operatorname{Cap}_2(E_k)\ge c
  \qquad(k\ge k_1).
\]
W szczególności szereg Wienera jest rozbieżny, a
$x_0$ jest regularnym punktem brzegowym, jeśli stosujemy
przywołane Twierdzenie~\ref{thm:pde19-wiener}.
\end{stwierdzenie}
\begin{proof}
Wybierzmy $\theta_0\in\Sigma$ i $\delta>0$ tak małe,
aby kierunki z kuli sferycznej wokół $\theta_0$ o promieniu
$2\delta$ należały do $\Sigma$, przy czym $\delta<1/16$.
Dla $h=2^{-k}<r_0$ weźmy
$z_k=x_0+3h\theta_0/4$. Możemy dodatkowo zmniejszyć
$\delta$, tak aby każdy kierunek o odległości euklidesowej
od $\theta_0$ nie większej niż $3\delta$ należał do $\Sigma$.
Jeśli $x\in\overline{B_{\delta h}(z_k)}$, to
$(3/4-\delta)h\le|x-x_0|\le(3/4+\delta)h$,
więc $h/2<|x-x_0|<h$. Dla $w=(x-x_0)/h$ mamy
$|w-(3/4)\theta_0|\le\delta$, a z nierówności trójkąta
\[
 \left|\frac{w}{|w|}-\theta_0\right|
 \le\frac{\bigl||w|-3/4\bigr|+|w-(3/4)\theta_0|}{3/4}
 \le\frac83\delta<3\delta.
\]
Kierunek należy zatem do $\Sigma$, a promień jest
mniejszy niż $r_0$. To dowodzi zawierania
$\overline{B_{\delta h}(z_k)}\subset E_k$.
Monotoniczność pojemności bezwzględnej daje
$\operatorname{Cap}_2(E_k)\ge
\operatorname{Cap}_2(\overline{B_{\delta h}(z_k)})$.
Przesunięcie nie zmienia pojemności, a z obliczenia energii
radialnej na pierścieniu i przejścia z jego promieniem
zewnętrznym do nieskończoności (dla $d\ge3$) mamy
$\operatorname{Cap}_2(\overline{B_\rho})
=(d-2)|S^{d-1}|\rho^{d-2}$.
Stąd dolne ograniczenie jest równe
$(d-2)|S^{d-1}|\delta^{d-2}h^{d-2}$, co po
pomnożeniu przez $h^{-(d-2)}$ dowodzi tezy.
\end{proof}

\begin{przyklad}[Przekłuta kula a kryterium Wienera]
\label{ex:pde19-wiener-puncture}
Dla $d\ge3$ oraz $\Omega=B_1(0)\setminus\{0\}$
w punkcie $x_0=0$ wszystkie dostatecznie małe pierścienie
$E_k$ są puste, gdyż dopełnienie obszaru zawiera
wewnątrz $B_1$ jedynie punkt $0$, niewchodzący do
żadnego z tych pierścieni. Suma z kryterium Wienera wynosi
zero, zatem $x_0$ jest nieregularny. Wynik zgadza się z
Przykładem~\ref{ex:pde19-irregular}.
\end{przyklad}



\begin{lemat}[Usuwalność ograniczonej osobliwości harmonicznej w wymiarze dwa]
\label{lem:pde-v5-removable-bounded-2d}
Niech $R>0$ i niech $h$ będzie rzeczywistą funkcją harmoniczną
na $B_R(0)\setminus\{0\}\subset\R^2$. Jeżeli $h$ jest ograniczona
w otoczeniu zera, to ma jednoznaczne przedłużenie harmoniczne
na całą kulę $B_R(0)$.
\end{lemat}
\begin{proof}
Wybierzmy $0<r_0<R$ i zapiszmy $x=(r\cos\theta,r\sin\theta)$.
Dla $0<r<r_0$ oznaczmy przez
\[
 c_n(r)=\frac1{2\pi}\int_{0}^{2\pi}
 h(r\cos\theta,r\sin\theta)e^{-in\theta}
 \,\mathrm d\lambda_1(\theta),\qquad n\in\mathbb Z,
\]
zespolone współczynniki Fouriera ograniczonej funkcji
$\theta\mapsto h(r\cos\theta,r\sin\theta)$.
Na każdym zwartym podpierścieniu możemy różniczkować
dwukrotnie pod całką, ponieważ pochodne $h$ są ciągłe.
Mnożąc równanie
$h_{rr}+r^{-1}h_r+r^{-2}h_{\theta\theta}=0$
przez $e^{-in\theta}$ i całkując po kącie,
dwukrotnie całkujemy przez części (wyrazy brzegowe
znikają przez okresowość). Dostajemy
$r^2c_n''+rc_n'-n^2c_n=0$. Podstawienie $r=e^t$
sprowadza je do równania o stałych współczynnikach
$\frac{\mathrm d^2}{\mathrm dt^2}c_n(e^t)-n^2c_n(e^t)=0$. Dla $n=0$ zatem
$c_0(r)=A_0+B_0\log r$, a dla $n\ne0$
$c_n(r)=A_nr^{|n|}+B_nr^{-|n|}$.
Ograniczoność $h$ oznacza $|c_n(r)|\le\|h\|_\infty$
dla dostatecznie małych $r$, więc wszystkie współczynniki
$B_n$ muszą znikać. Niech $H$ będzie harmonicznym
przedłużeniem Poissona funkcji $h|_{\partial B_{r_0}}$
do $B_{r_0}(0)$ (sekcja~\ref{sec:pde14-poisson}).
Wzór Poissona na dysku wykorzystuje jądro
$(1-q^2)/(1-2q\cos\alpha+q^2)$, gdzie $q=r/r_0<1$.
Z dwóch szeregów geometrycznych wynika rozwinięcie
tego jądra jako $\sum_{n\in\mathbb Z}q^{|n|}e^{in\alpha}$,
jednostajnie i bezwzględnie zbieżne dla ustalonego $q$.
Całkując szereg przeciwko danym brzegowym, otrzymujemy
współczynniki $(r/r_0)^{|n|}c_n(r_0)$; dla $n=0$
są stałe.
Są zatem identyczne z $c_n(r)$ dla $0<r<r_0$.
Z jednoznaczności szeregów Fouriera dla funkcji ciągłych
na okręgu wynika $H=h$ w przekłutej kuli $B_{r_0}$.
Funkcja otrzymana przez dołączenie $H(0)$ jest więc
harmonicznym przedłużeniem $h$ na całą $B_R$.
Jednoznaczność wynika z ciągłości.
\end{proof}

\begin{stwierdzenie}[Nieregularność izolowanego punktu brzegu w wymiarze dwa]
\label{prop:pde-v5-punctured-disk}
Niech $\Omega=B_1(0)\setminus\{0\}\subset\R^2$.
Punkt $0\in\partial\Omega$ nie jest regularny dla
klasycznego zagadnienia Dirichleta w sensie Perrona.
Dokładniej, dla danych brzegowych
$f(0)=1$ oraz $f|_{\partial B_1}=0$ rozwiązanie
Perrona jest identycznie równe zeru w $\Omega$.
\end{stwierdzenie}
\begin{proof}
Funkcja $f$ jest ciągła na rozłącznej sumie
$\partial\Omega=\{0\}\cup\partial B_1$.
Z zasady porównania ze stałymi wynika, że rozwiązanie
Perrona $u$ jest ograniczone przez $0\le u\le1$;
standardowa konstrukcja Perrona zapewnia jego harmoniczność
wewnątrz $\Omega$.
Z lematu~\ref{lem:pde-v5-removable-bounded-2d} funkcja $u$
przedłuża się harmonicznie przez zero. Każdy punkt
zewnętrznej sfery $\partial B_1$ ma barierę pochodzącą
od kuli zewnętrznej, więc $u$ dąży tam do wartości $0$
(przywołujemy tu standardową własność rozwiązania Perrona
w punktach z barierą).
Zasada maksimum dla funkcji harmonicznych w $B_1$
implikuje $u\equiv0$. Stąd granica $u(x)$ przy $x\to0$
wynosi $0$, a nie $f(0)=1$, co dowodzi nieregularności.
Argument nie wymaga pełnego logarytmicznego kryterium Wienera.
\end{proof}

\begin{przyklad}[Rodzina punktów regularnych i nieregularnych
sterowana szeregiem]
\label{ex:pde19-wiener-balls}
Niech $d\ge3$, $e_1=(1,0,\ldots,0)$, $h_k=2^{-k}$,
$k\ge2$, oraz $0<\delta_k\le1/16$. Określmy
$z_k=\tfrac34h_ke_1$, $K_k=\overline{B_{\delta_kh_k}(z_k)}$
i
\[
 \Omega=B_1(0)\setminus\biggl(\{0\}\cup
 \bigcup_{k=2}^\infty K_k\biggr).
\]
Wówczas $\Omega$ jest ograniczonym, otwartym i spójnym
obszarem, a $0\in\partial\Omega$. Kryterium
Wienera~\ref{thm:pde19-wiener} daje dokładnie
\[
 0\text{ jest regularny}\quad\Longleftrightarrow\quad
 \sum_{k=2}^\infty\delta_k^{d-2}=+\infty.
\]
W szczególności, dla $\delta_k=\frac1{16}k^{-a}$,
$a\ge0$, punkt $0$ jest regularny wtedy i tylko wtedy,
gdy $a(d-2)\le1$.
\end{przyklad}
\begin{proof}
Kule $K_k$ są parami rozłączne. Ich środki leżą na
promieniach $3h_k/4$, a promienie nie przekraczają
$h_k/16$. Dlatego $K_k$ zawiera się w otwartym
pierścieniu $h_k/2<|x|<h_k$. Pozostałe kule nie przecinają
tego pierścienia; zbiór $E_k$ z definicji kryterium Wienera
jest więc równy $K_k$. Usunięty zbiór jest domknięty
w $B_1$, ponieważ jego jedynym punktem skupienia poza
kulami jest $0$; dołączono ten punkt. Ponieważ punkty $-te_1$ dla $0<t<1$ należą do
$\Omega$, punkt $0$ jest rzeczywiście punktem brzegu.
Wykażemy łukową spójność. Dla dowolnych $x,y\in\Omega$
wybieramy $m$ tak duże, by $h_m<\min(|x|,|y|)$.
Kule $K_k$ dla $k\ge m$ leżą wewnątrz $B_{h_m}$,
a kule dla $k<m$ leżą poza $\overline{B_{h_m}}$.
Wszystkie domknięte kule z tej skończonej rodziny
(wraz z $\overline{B_{h_m}}$) są parami rozłączne,
leżą wewnątrz $B_1$ i nie zawierają $x,y$.
Możemy nieznacznie zwiększyć każdy z ich promieni,
zachowując te własności. Odcinek od $x$ do $y$ leży
w $B_1$. Każdy jego fragment przechodzący przez jedną
z powiększonych kul zastępujemy łukiem na jej sferze
(która jest łukowo spójna dla $d\ge2$).
Sfery są rozłączne z pozostałymi powiększonymi kulami,
więc po skończenie wielu zastąpieniach otrzymujemy
drogę w $\Omega$. To dowodzi spójności.

Dla $d\ge3$ bezwzględna pojemność kuli ma wzór
$\operatorname{Cap}_2(K_k)=(d-2)|S^{d-1}|
(\delta_kh_k)^{d-2}$ (Twierdzenie~\ref{thm:pde19-balls}
z promieniem zewnętrznym dążącym do nieskończoności).
Stąd
\[
 2^{k(d-2)}\operatorname{Cap}_2(E_k)
 =(d-2)|S^{d-1}|\delta_k^{d-2}.
\]
Szereg Wienera jest zatem rozbieżny dokładnie wtedy,
gdy rozbieżny jest wskazany szereg liczbowy. Ostatnia
równoważność wynika z kryterium zbieżności szeregu
$\sum_{k\ge2}k^{-a(d-2)}$.
\end{proof}

\begin{uwaga}[Granica samodzielności argumentu]
Geometria kul, obliczenie pojemności i redukcja do szeregu
liczbowego zostały wyprowadzone powyżej. Równoważność
między rozbieżnością szeregu a regularnością punktu brzegu
jest zastosowaniem \emph{przywołanego}, a nie udowodnionego
tutaj kryterium Wienera.
\end{uwaga}

\section{Związki z dalszą regularnością i podsumowanie}
\label{sec:pde19-summary}
Pojemność rozróżnia zbiory, na których w naturalny sposób
może koncentrować się energia rozwiązań, od tych, które
są dla przestrzeni Sobolewa niewidoczne. Wykazaliśmy
trzy konkretne mechanizmy: znikanie pojemności punktów
w wymiarach $d\ge2$, przedłużanie równań dywergencyjnych
przez zbiory pojemności zero dla rozwiązań $H^1$ oraz
istnienie quasi-ciągłych reprezentantów funkcji Sobolewa.
Kryterium Wienera opisuje natomiast regularność warunku
brzegowego za pomocą rozłożenia pojemności na skalach.

W późniejszej teorii jedyności kontynuacji i analizie
mikrolokalnej pojawią się inne sposoby mierzenia wielkości
zbiorów wyjątkowych; nie są one równoważne pojemności
w dowolnych klasach współczynników.


% Rozdzial 20 -- samodzielny fragment do wlaczenia przed \appendix.
